\documentclass[aspectratio=169]{beamer}

\usetheme{Madrid}
\usecolortheme{seahorse}
\usepackage{graphicx}
\usepackage{booktabs}
\usepackage{siunitx}

\title{AM CVn: Testing a Superhump Interpretation}
\author{TODO: Team Members}
\date{PHYSICS 100 Final Project}

\begin{document}

\begin{frame}
  \titlepage
\end{frame}

\begin{frame}{Scientific Goal}
  \begin{itemize}
    \item AM CVn is an ultracompact helium-accreting binary.
    \item Literature periods: orbital \(1028.7322\) s, positive superhump \(1051.2\) s, strong harmonic near \(525.6\) s.
    \item Question: can our 2026 May 20 photometry tighten those constraints?
    \item Rule: separate free-period constraints from literature-conditioned fits.
  \end{itemize}
\end{frame}

\begin{frame}{Observations}
  \begin{columns}
    \begin{column}{0.50\textwidth}
      \begin{itemize}
        \item 140 V-band frames.
        \item 35, 40, and 45 s exposures.
        \item 138 good points after rejecting two artifact frames.
        \item Baseline: 111.20 min.
        \item Covers 6.35 cycles of a 1051.2 s superhump.
      \end{itemize}
    \end{column}
    \begin{column}{0.50\textwidth}
      \includegraphics[width=\textwidth]{analysis/project/plots/AM_CVn_sequence_X42421BZ_light_curve.png}
    \end{column}
  \end{columns}
\end{frame}

\begin{frame}{Reduction}
  \begin{itemize}
    \item Bias, dark, and star-masked flat-field calibration.
    \item AAVSO sequence X42421BZ used for target and comparison-star geometry.
    \item Fluxes converted to rates before combining different exposure times.
    \item Comparison stars normalized by their own median before forming the ensemble.
    \item Robust target scatter: 1.51 percent; median formal error: 0.90 percent.
  \end{itemize}
\end{frame}

\begin{frame}{Model Ladder}
  \begin{table}
    \small
    \begin{tabular}{lrrr}
      \toprule
      Model & Period & \(\Delta\)BIC & CV RMS \\
      \midrule
      Free double wave & 1055.21 s & 0.00 & 1.52\% \\
      Beat-constrained disk & 1055.22 s & 0.00 & 1.52\% \\
      1051.2 s + harmonic & fixed & 0.55 & 1.53\% \\
      Free single sine & 525.06 s & 37.44 & 1.57\% \\
      Orbital + superhumps & fixed & 45.55 & 1.65\% \\
      Superhump only & 1051.2 s & 57.67 & 1.62\% \\
      Orbital only & 1028.73 s & 66.18 & 1.65\% \\
      \bottomrule
    \end{tabular}
  \end{table}
\end{frame}

\begin{frame}{Best Physical Interpretation}
  \centering
  \includegraphics[height=0.78\textheight]{analysis/project/plots/AM_CVn_physical_model_ladder.png}
\end{frame}

\begin{frame}{Constraints}
  \begin{columns}
    \begin{column}{0.47\textwidth}
      \begin{itemize}
        \item Free single sine: \(525.06 \pm 4.27\) s.
        \item Free double wave: \(1055.21 \pm 8.58\) s.
        \item Conditional superhump family: \(1055.22 \pm 8.64\) s.
        \item These are consistent with the 1051.2 s superhump family.
      \end{itemize}
    \end{column}
    \begin{column}{0.53\textwidth}
      \includegraphics[width=\textwidth]{analysis/project/plots/AM_CVn_physical_literature_constraints.png}
    \end{column}
  \end{columns}
\end{frame}

\begin{frame}{Validation}
  \small
  \begin{columns}
    \begin{column}{0.50\textwidth}
      \begin{itemize}
        \item Multi-harmonic Lomb-Scargle recovers the double-wave fundamental.
        \item Beat-constrained model implies \(P_{\rm prec}\approx 11.39\) hr, inferred rather than directly observed.
        \item Leave-one-comparison-star-out period span: 5.2 s.
        \item Outlier clipping period shift: 0.02 s.
      \end{itemize}
    \end{column}
    \begin{column}{0.50\textwidth}
      \begin{itemize}
        \item Residual periodogram has no significant remaining peak.
        \item Systematics terms lower RMS to 1.41\%, but worsen BIC/CV.
        \item Robust losses shift the answer to 1057--1059 s and worsen RMS.
        \item First/second halves differ by 22.8 s.
      \end{itemize}
    \end{column}
  \end{columns}
\end{frame}

\begin{frame}{Could It Be Something Else?}
  \begin{itemize}
    \item Broadband photometry alone cannot prove the object class.
    \item Literature spectroscopy shows helium accretion signatures and an S-wave orbital feature.
    \item A non-binary pulsator would need a spectrum without accretion features and a different multi-period pattern.
    \item The decisive test is time-resolved spectroscopy, not another photometric periodogram.
  \end{itemize}
\end{frame}

\begin{frame}{Conclusion}
  \begin{itemize}
    \item We detect AM CVn short-timescale variability.
    \item The 525 s single-sine signal is best interpreted as a harmonic.
    \item The double-wave and 1051.2 s superhump-template fits support the known superhump/harmonic structure.
    \item Detrending, robust likelihoods, and comparison subsets do not improve the constraint.
    \item We do not tighten the literature periods with this dataset.
    \item A longer multi-night baseline is the path to a stronger constraint.
  \end{itemize}
\end{frame}

\end{document}
